\documentclass[12pt]{article}
\usepackage[T1]{fontenc}
\usepackage[utf8]{inputenc}
\usepackage{lmodern}
\usepackage{textcomp}
\usepackage{enumitem}
\usepackage{geometry}
\geometry{margin=1in}
\begin{document}
\begin{center}
Answer Key - Biology - Graduate Level Assessment 15 \
\small (Answers are ordered exactly as in the question paper.)
\end{center}

\section*{Multiple Choice}
\begin{enumerate}[label=\arabic*.]
\item \textbf{Answer:} A
\\ \textit{Options:} A) Golgi apparatus and endoplasmic reticulum; B) Lysosomes and endosomes; C) Plasma membrane and cytoskeleton; D) DNA and mRNA; E) Proteins and lipids
\\ \textit{Note:} The Golgi apparatus and endoplasmic reticulum can be separated by low-speed centrifugation due to their differences in size and density, allowing for their distinct organelles to pellet at different rates during centrifugation.
\item \textbf{Answer:} C
\\ \textit{Options:} A) They can reproduce asexually.; B) They can produce viable gametes through mitosis.; C) They have a relative evolutionary fitness of zero.; D) Their offspring are always genetically identical to the parents.; E) They cannot produce offspring due to lack of genetic variation.
\\ \textit{Note:} Mules are sterile hybrids resulting from the crossbreeding of a horse and a donkey, which prevents them from producing viable gametes and thus leads to a relative evolutionary fitness of zero, as they cannot pass on their genes to future generations.
\item \textbf{Answer:} A
\\ \textit{Options:} A) Acetylcholine will not be degraded in the synaptic cleft.; B) Acetylcholine will not bind to receptor proteins on the postsynaptic membrane.; C) Acetylcholine will not be produced in the presynaptic membrane.; D) Acetylcholine will be excessively produced in the synaptic cleft.; E) Acetylcholine will not be transported to the synaptic cleft.
\\ \textit{Note:} The correct answer is based on the fact that acetylcholine is normally degraded by the enzyme acetylcholinesterase in the synaptic cleft, so if this enzyme is denatured by the chemical agent, acetylcholine will remain in the cleft and not be broken down.
\item \textbf{Answer:} A
\\ \textit{Options:} A) Both species benefit equally; B) One species benefits, and the other is significantly harmed; C) One species harms the other without any benefit; D) Both species harm each other; E) The species exchange benefits, but one benefits more than the other
\\ \textit{Note:} In a commensalistic relationship, one species benefits while the other is neither helped nor harmed, which contradicts the assertion that both species benefit equally, highlighting that the correct answer should reflect the unequal nature of benefits in such interactions.
\item \textbf{Answer:} A
\\ \textit{Options:} A) 4; B) 16; C) 32; D) 64; E) 192
\\ \textit{Note:} Meiosis reduces the chromosome number by half, so an organism with a diploid number (2 n) of 96 would produce gametes with a haploid number (n) of 48 chromosomes, and since the question seems to imply a misunderstanding, the answer of 4 appears to be a misinterpretation; the correct answer should reflect the halving of 96 to 48, indicating clarification is needed.
\end{enumerate}

\section*{True / False}
\begin{enumerate}[label=\arabic*.]
\setcounter{enumi}{5}
\item \textbf{Answer:} False
\\ \textit{Options:} A) True; B) False
\\ \textit{Note:} The incidence of CPM is not increased in IVF/ICSI pregnancies compared with spontaneous conceptions. CPM probably does not account for the adverse perinatal outcomes following IVF/ICSI.
\item \textbf{Answer:} False
\\ \textit{Options:} A) True; B) False
\\ \textit{Note:} The effect on venous pressures caused by the change in patient positioning alone during liver surgery does not affect the risk of venous air embolism.
\item \textbf{Answer:} True
\\ \textit{Options:} A) True; B) False
\\ \textit{Note:} The incidence of post-tonsillectomy late haemorrhage in our study population was 1.78\%. A statistically significant difference was found between night-time and day-time haemorrhages. Even though no significant distribution of haemorrhages per hour was observed, we underline that we recorded 32 (54.2\%) events in 2 periods of the day: from 10 p.m. to 1 a.m. and from 6 to 9 a.m.
\item \textbf{Answer:} True
\\ \textit{Options:} A) True; B) False
\\ \textit{Note:} It is important to continue FA supplementation over the long term in patients on methotrexate and FA in order to prevent them discontinuing treatment because of mouth ulcers or nausea and vomiting. Our data suggest that FA supplementation is also helpful in preventing neutropenia, with very little loss of efficacy of methotrexate.
\item \textbf{Answer:} True
\\ \textit{Options:} A) True; B) False
\\ \textit{Note:} The presence of a high signal intensity change on T$_{2}$-weighted MRI and the pyramidal tract sign can be used as prognostic factors for patients with CSA.
\end{enumerate}

\section*{Long Form Response}
\begin{enumerate}[label=\arabic*.]
\setcounter{enumi}{10}
\item \textbf{Answer:} The lower frequency of the sickle-cell anemia gene among African-Americans compared to populations in Africa can be attributed to a combination of environmental and social factors, particularly the increased mortality associated with sickle-cell disease in the United States. In Africa, the prevalence of the sickle-cell trait is maintained due to its protective advantage against malaria, which is a significant selective pressure in tropical regions. However, in the U.S., individuals with sickle-cell disease often face severe health complications and lower life expectancy due to limited access to healthcare, socioeconomic disparities, and a lack of awareness about the disease. These factors contribute to a higher mortality rate among affected individuals, leading to a reduced frequency of the gene in African-American populations. Consequently, while the sickle-cell trait offers a survival benefit in certain environmental contexts, the social and healthcare dynamics in the U.S. negate this advantage, resulting in lower allele frequencies.
\\ \textit{Note:} correct because it highlights how the protective advantage of the sickle-cell trait against malaria in Africa contrasts with the higher mortality associated with sickle-cell disease in the U.S., leading to a lower frequency of the gene among African-Americans due to differing environmental selective pressures and social health outcomes.
\item \textbf{Answer:} All living organisms share fundamental nutritional requirements, including the need for energy, carbon, nitrogen, vitamins, and minerals to sustain metabolic processes and growth. These requirements manifest differently among various groups; for instance, phototrophs harness light energy to convert inorganic substances into organic matter, while chemotrophs obtain energy from chemical compounds, often through oxidation processes. Autotrophs, such as plants, produce their own food using carbon dioxide and sunlight or inorganic compounds, whereas heterotrophs rely on consuming other organisms for energy and organic carbon. While all organisms must satisfy these basic needs, the mechanisms and sources from which they derive their nutrients highlight significant ecological and evolutionary differences. Notably, only animals and plants have distinct nutritional requirements that categorize them into these groups, reflecting their adaptations to specific environments and ecological roles.
\\ \textit{Note:} The answer correctly identifies the universal nutritional requirements of energy, carbon, nitrogen, vitamins, and minerals while effectively differentiating the methods of energy acquisition and organic matter production among phototrophs, chemotrophs, autotrophs, and heterotrophs.
\end{enumerate}

\vfill
\noindent\textit{End of Answer Key}
\end{document}